\documentclass[11pt,a4paper]{article}

\usepackage[margin=1in]{geometry}
\usepackage[T1]{fontenc}
\usepackage[utf8]{inputenc}
\usepackage{hyperref}
\usepackage{enumitem}
\usepackage{titlesec}
\usepackage{xparse}

\setlength{\parindent}{0pt}
\setlength{\parskip}{6pt}

\titleformat{\section}{\large\bfseries}{}{0pt}{}
\titleformat{\subsection}{\normalsize\bfseries}{}{0pt}{}


\NewDocumentCommand{\meet}{m o m o}{%
\subsection*{#1}
% \textbf{Progress:} #2
\IfNoValueF{#2}{\textbf{Progress:} #2}

\textbf{Next Steps:} #3

\IfNoValueF{#4}{\textbf{Notes:} #4}
}


\begin{document}

\section*{Meeting Log}


\textbf{Project:} Automated Parallel Goal-Based Adaptivity\\
\textbf{Student:} Yue Wu\\
\textbf{Supervisor:} Prof. Patrick Farrell\\



\meet
{Meeting on 2026-08-24}
[Completed the first thesis draft, and tested 4 different problems.]
{\begin{itemize}
\item Add another goal functional test for p-Laplace and add axis.
\item Abstract: perhaps the most important part.
\item Try to solve Navier-Stokes from Taylor-hood.
\end{itemize}}
\vspace{1.5em}



\meet
{Meeting on 2026-08-19}
[Revise the 1D example, start to working on BBM, and organise the code.]
{Thesis writing and more numerical tests. Try other source term for heat.}
\vspace{1.5em}



\meet
{Meeting on 2026-08-12}
[Compared manually derived and automated 2d heat results and applied the automated solver on a 1d advection-diffusion equation ]
{\begin{itemize}
\item Revise the refinement for the 1d equation.
\item 8.24 submit the draft!
\end{itemize}}
[\begin{itemize}
\item The refinement strategy - how can we explain it? Convergence results from dorfler for stationary case..
\item Observe the meshes for heat - can’t tell the difference. Should try different analytical solutions.
\item To choose the tolerance - try 10, 20,... iterations.
\item For experiments, check the Irksome demo and paper. e.g., BBM equation, superposition of solitons, incompressible Navier-Stokes...
\end{itemize}]
\vspace{1.5em}



\meet
{Meeting on 2026-07-24 and 2026-08-03}
[Finished the draft of the first two chapters for MSc thesis. Written down the proposal for automated solver.]
{Implement the bubble recovery localisation, and compare it with the previous heat equation results.
\begin{itemize}
\item Do the refinement (to different basis mesh) and see if they re broadly similar.
\item Observe how much extra localisation efficiency index is needed.
\item Observe how much DoFs are ended up with to achieve the same goal.
\item To make it concrete, we can prove how much the estimator would change.
\end{itemize}}
\vspace{1.5em}



\meet
{Meeting on 2026-07-17}
[Explored two different goal functionals and completed the implementation of four temporal enrichment strategies.]
{\begin{itemize}
\item Commit the latest code to Git.
\item Begin investigating the core mathematics of the general framework.
\item Code: examine the order of temporal and spatial refinement within each adaptive iteration.
\item Start drafting the background chapter of the thesis.
\end{itemize}}
[Use $p$-enrichment in the error estimator and $h$-refinement in the adaptive procedure. Once the adaptive process has been automated, consider testing the DWR method on an ODE and a one-dimensional wave problem. This would allow the temporal and spatial refinement patterns to be visualised together in a two-dimensional plot.]
\vspace{1.5em}



\meet
{Meeting on 2026-07-06}
[Completed two experiments to compare manual derivation results and Irksome results. However, for DG1, the convergence rate did not exhibit a well-behaved pattern as the mesh was refined or the time step was reduced.]
{\begin{itemize}
    \item Compare the results with the bisect approach using ParaView:
    \begin{itemize}
        \item Check whether the numerical solutions differ visually.
        \item Investigate whether the space-time partitioning needs to be adjusted.
    \end{itemize}
    \item Consider how to implement localised adaptive refinement.
\end{itemize}}
[Keep the user interface in mind during development. And codex (CLI) is helpful.]
\vspace{1.5em}


\meet
{Meeting on 2026-06-25}
[Complete a preliminary adaptive workflow and assess the error estimation strategy for the heat equation.]
{\begin{itemize}
    \item Use Irksome to compute the adjoint automatically.
    \item Check whether the global error estimate coincides.
    \item Keep the refinement pattern unchanged for the comparison.
\end{itemize}
}
[Error estimate from DG1.]
\vspace{1.5em}


\meet
{Meeting on 2026-06-18}
[Extending the bubble projection idea from the stationary case to the space-time setting. Yet the effective indices do not approach 1; instead, they keep increasing, which suggests a problem in the enrichment treatment.]
{Debug the space-time implementation by starting with a 1D space-time boundary value problem.}
\vspace{1.5em}



\meet
{Meeting on 2026-06-11}
[Reviewing the literature and comparing different strategies.]
{Implement the complete workflow using the heat equation.}
[\begin{itemize}
    \item For now, do not consider iteration error (keep everything simple at the manual derivation stage).
    \item Interpolation-based local higher-order approximations was previously constructed due to computational performance constraints. Now enrichment can be used directly at relatively low cost.
    \item Operations such as integration by parts, which require symbolic manipulation, should be avoided in the automated workflow.
    \item[*] In the previous stationary case, the geometry error issue was resolved by formulating the problem at the PDE level as a mixed formulation, including the mapping of the domain. (Similar to the eigenvalue problem, the key is to start from a good variational formulation.)
\end{itemize}]
\vspace{1.5em}



\meet
{Meeting on 2026-06-03}
[Preparation for the presentation.]
{Presentation and slide design strategies:
\begin{itemize}
    \item Use motivating examples whenever possible.
    \item Use pauses and overlays to improve the flow of the presentation.
    \item For theory-heavy slides, consider allocating a small column (e.g. 20\%) for photographs of key researchers.
    \item Use colour coding to distinguish positive and negative results.
\end{itemize}
}
[\begin{itemize}
    \item Clarified the central role of the dual problem and localisation of the error representation.
    \item The enrichment strategy should use $p$-enrichment for space.
    \item Discussed theoretical concepts such as the saturation assumption and enrichment indices, which are not entirely satisfactory from a theoretical perspective. DWR methods often perform significantly better in practice than their current theoretical guarantees suggest, similar to GMRES and Newton's method.
\end{itemize}]
\vspace{1.5em}



\meet
{Meeting on 2026-05-28}
[Code implementation for global error estimation.]
{Preparation for the dissertation presentation.}
[\begin{itemize}
    \item Challenges in the stationary case: derivation of the adjoint problem, and obtaining the strong form through localization.
    \item Main project focus: localisation of the error representation, and adaptive balancing of spatial and temporal errors.
    \item Time discretization in the code: DG(0) versus backward Euler.
    \item Verification of the adjoint implementation using https://www.dolfin-adjoint.org/en/latest/. (Correctness can be assessed through derivative convergence tests, which should exhibit better than first-order convergence.)
    \item Investigated whether temporal enrichment may also be required when using enriched spatial spaces.
\end{itemize}
]
\vspace{1.5em}



\meet
{Meeting on 2026-05-20}
[Basic understanding of the framework: 
weak formulation $\rightarrow$ space-time finite elements $\rightarrow$ goal functionals $\rightarrow$ DWR $\rightarrow$ error identity $\rightarrow$ error estimation in enriched spaces]
{Code implementation: 2D heat equation example.
Compare the error in $J$ with the known exact solution.}
[choice of goal functional (time regularity)]
\vspace{1.5em}



\meet
{Meeting on 2026-05-14}
{Working on the 2D heat equation example.
}
[DG0 in time: backward Euler]
% [
% \begin{itemize}
%     \item DG0 in time.
% \end{itemize}
% ]



\end{document}